\documentclass[11pt,a4paper]{article}

% Shared front-matter data. Edit here, not in the report bodies.

\newcommand{\coursename}{Blockchain Technology (BITS F452)}
\newcommand{\groupnumber}{6}
\newcommand{\projectname}{DVote}

% Columns are sized to their content so the rules always span the full table.
\newcommand{\groupmemberstable}{%
  \small
  \begin{tabular}{@{}lll@{}}
    \toprule
    \textbf{ID No.} & \textbf{Name} & \textbf{Email} \\
    \midrule
    2024A7PS0148H & Varad Kailas Thombre & \texttt{f20240148@hyderabad.bits-pilani.ac.in} \\
    2024A7PS0121H & Naman Singhal        & \texttt{f20240121@hyderabad.bits-pilani.ac.in} \\
    2024ADPS0190H & Shivam Goyal         & \texttt{f20240190@hyderabad.bits-pilani.ac.in} \\
    2024A7PS0101H & Rishit Agarwal       & \texttt{f20240101@hyderabad.bits-pilani.ac.in} \\
    2024A7PS0113H & Vihaan Kashyap       & \texttt{f20240113@hyderabad.bits-pilani.ac.in} \\
    2024A7PS0041H & Palak Agarwal        & \texttt{f20240041@hyderabad.bits-pilani.ac.in} \\
    2023B4A70727H & Chaitanya Keyal      & \texttt{f20230727@hyderabad.bits-pilani.ac.in} \\
    2024B4A70830H & Wamika Sehgal        & \texttt{f20240830@hyderabad.bits-pilani.ac.in} \\
    \bottomrule
  \end{tabular}%
}


\newcommand{\reporttitle}{\projectname: A Decentralized Voting System}
\newcommand{\reportsubtitle}{Evaluation 2: Final Report}
\newcommand{\submissiondate}{November 2026}

% Shared preamble for both evaluation reports.
% Compile with: latexmk -pdf eval1.tex   (see Makefile)

\usepackage[T1]{fontenc}
\usepackage[utf8]{inputenc}
\usepackage{lmodern}
% The .docx template uses Arial throughout; Helvetica is its metric-compatible
% LaTeX equivalent.
\usepackage{helvet}
\renewcommand{\familydefault}{\sfdefault}
% Margins and header spacing match the template: 0.5in all round, with room for
% the institute logo that sits in the header of every page.
\usepackage[a4paper,margin=0.5in,headheight=88pt,headsep=14pt,includehead]{geometry}
\usepackage{microtype}

\usepackage{amsmath,amssymb}
\usepackage{graphicx}
\usepackage{booktabs}
\usepackage{longtable}
\usepackage{array}
\usepackage{ragged2e}
\usepackage{enumitem}
\usepackage{xcolor}
\usepackage{caption}
\usepackage{fancyhdr}
\usepackage{titlesec}
\usepackage{listings}
\usepackage{csquotes}
\usepackage{tikz}
\usetikzlibrary{arrows.meta,positioning,shapes.geometric,fit,backgrounds,calc}
% Graphs are drawn from the committed CSVs in bench/results, never by hand.
\usepackage{pgfplots}
\usepackage{pgfplotstable}
\pgfplotsset{compat=1.18}
\usepgfplotslibrary{groupplots}

\usepackage[backend=biber,style=numeric,sorting=none]{biblatex}
\addbibresource{references.bib}

\usepackage[hidelinks,breaklinks]{hyperref}

% Keep the PDF free of metadata: no title, author, creator or producer
% strings, no timestamps, no TeX version banner and no random trailer ID.
\hypersetup{pdftitle={}, pdfauthor={}, pdfsubject={}, pdfkeywords={},
            pdfcreator={}, pdfproducer={}}
\pdfinfoomitdate=1
\pdftrailerid{}
\pdfsuppressptexinfo=-1

\graphicspath{{figures/}}

% ---------------------------------------------------------------------------
% Draft placeholders
%
% \stub{...} marks unwritten content. It renders as a visible red box, so an
% unfinished section is obvious on the page. Delete the \stub call as you write
% the prose that replaces it; when none are left, the report is done.
% ---------------------------------------------------------------------------
\newcommand{\stub}[1]{%
  \par\noindent
  \begingroup
  \color{red!65!black}\small\itshape
  \fbox{\parbox{\dimexpr\linewidth-2\fboxsep-2\fboxrule}{TODO: #1}}%
  \endgroup
  \par\medskip
}

% Inline marker for a single value still to be filled in.
\newcommand{\note}[1]{\textcolor{blue!55!black}{[#1]}}

% ---------------------------------------------------------------------------
% Layout
% ---------------------------------------------------------------------------
% Every page carries the institute logo in the header, as the template does.
\newcommand{\institutelogo}{\includegraphics[width=5.81in]{bits-logo}}

\pagestyle{fancy}
\fancyhf{}
\fancyhead[C]{\institutelogo}
\fancyfoot[C]{\thepage}
\renewcommand{\headrulewidth}{0pt}

% Title page: logo header, no page number.
\fancypagestyle{titlestyle}{%
  \fancyhf{}%
  \fancyhead[C]{\institutelogo}%
  \renewcommand{\headrulewidth}{0pt}%
}

% The plain style is used on the first page of the document and by \chapter-like
% breaks; give it the same header so no page is missing the logo.
\fancypagestyle{plain}{%
  \fancyhf{}%
  \fancyhead[C]{\institutelogo}%
  \fancyfoot[C]{\thepage}%
  \renewcommand{\headrulewidth}{0pt}%
}

\titleformat{\section}{\large\bfseries}{\thesection.}{0.6em}{}
\titleformat{\subsection}{\normalsize\bfseries}{\thesubsection}{0.6em}{}

\setlist[itemize]{topsep=2pt,itemsep=2pt,parsep=0pt,leftmargin=1.4em}
\setlist[enumerate]{topsep=2pt,itemsep=2pt,parsep=0pt,leftmargin=1.6em}

\captionsetup{font=small,labelfont=bf}

% A page holding only floats is centred vertically by default, which strands a
% figure in the middle of an otherwise empty page. Push those pages to the top.
\makeatletter
\setlength{\@fptop}{0pt}
\makeatother

% Fixed-width column types for tables that need wrapping.
\newcolumntype{L}[1]{>{\RaggedRight\arraybackslash}p{#1}}
\newcolumntype{R}[1]{>{\RaggedLeft\arraybackslash}p{#1}}

\lstset{
  basicstyle=\ttfamily\footnotesize,
  breaklines=true,
  frame=single,
  framesep=4pt,
  rulecolor=\color{black!30},
  numbers=left,
  numberstyle=\tiny\color{black!50},
  keywordstyle=\bfseries,
  showstringspaces=false,
  captionpos=b,
}

% Solidity is not a built-in listings language; define enough to be readable.
\lstdefinelanguage{Solidity}{
  keywords={pragma,solidity,contract,library,interface,function,modifier,event,
      emit,struct,enum,mapping,memory,storage,calldata,public,private,
      internal,external,view,pure,payable,returns,return,require,revert,
      if,else,for,while,new,constructor,immutable,constant,is,using,
      address,uint256,bytes32,bool,string,error},
  keywordstyle=\bfseries,
  comment=[l]{//},
  morecomment=[s]{/*}{*/},
  commentstyle=\color{black!55}\itshape,
  string=[b]",
  stringstyle=\color{black!70},
  morestring=[b]',
}

% ---------------------------------------------------------------------------
% Title page
% ---------------------------------------------------------------------------
\newcommand{\maketitlepage}{%
  % The title page needs no hyperref anchor, and suppressing it avoids a
  % duplicate page.1 destination once front matter switches to roman numerals.
  \hypersetup{pageanchor=false}%
  \begin{titlepage}
    \thispagestyle{titlestyle}
    \centering
    \vspace*{1.5cm}

    {\Huge\bfseries \reporttitle\par}
    \vspace{0.8cm}
    {\Large \reportsubtitle\par}

    \vspace{1.5cm}
    {\large \coursename\par}
    \vspace{0.3cm}
    {\large \textbf{Group \groupnumber}\par}

    \vspace{1.2cm}
    {\large\bfseries Group Members\par}
    \vspace{0.4cm}
    \groupmemberstable

    \vfill
  \end{titlepage}
  % Front matter in roman numerals so the title page and the body do not both
  % claim page 1 (which produces duplicate hyperref anchors).
  \pagenumbering{roman}%
  \hypersetup{pageanchor=true}%
}

% Ends the front matter and starts the numbered body.
\newcommand{\startbody}{\clearpage\pagenumbering{arabic}}


\begin{document}

\maketitlepage

\tableofcontents
\startbody

% ===========================================================================
\section{Introduction}
% Sections 1-3 are carried over from the Evaluation 1 report and revised to
% reflect what was actually built. Do not rewrite them from scratch; do not
% paste them unchanged either; the design will have moved.
% ===========================================================================

\stub{Revise the Evaluation 1 introduction. State what was built, and add one
  paragraph summarising the headline results so a reader knows early what the
  report concludes.}

% ===========================================================================
\section{Domain and Problem Analysis}
% ===========================================================================

\stub{Carry over from Evaluation 1. Revise anything the implementation
  contradicted. If a problem turned out to be different in practice than it
  looked on paper, say so.}

% ===========================================================================
\section{Requirement Gathering}
% ===========================================================================

\stub{Carry over the actors and requirements from Evaluation 1, then add a short
  subsection recording which requirements changed during implementation and why.
  Requirements that were dropped are as informative as requirements that were met.}

% ===========================================================================
\section{Final System Architecture}
% Source material: docs/architecture.md, as built
% ===========================================================================

\subsection{Architecture as built}

\stub{Present the final architecture diagram and describe the deployed system.
  Where it differs from the proposal in Evaluation 1, state the difference and the
  reason for it.}

\begin{figure}[htbp]
  \centering
  \fbox{\parbox[c][5.5cm][c]{0.9\linewidth}{\centering\itshape
      Figure placeholder: final architecture.}}
  \caption{Deployed system architecture.}
  \label{fig:final-architecture}
\end{figure}

\subsection{Deployment}

\stub{Record the deployed contract addresses, the network, the block explorer
  links, and the commit hash the deployment was built from. A reader should be able
  to independently verify every claim in Section 6 from this information alone.}

\begin{table}[htbp]
  \centering
  \caption{Deployed contracts.}
  \label{tab:deployment}
  \small
  \begin{tabular}{@{}L{4.6cm}L{7.4cm}L{3.9cm}@{}}
    \toprule
    \textbf{Contract}                & \textbf{Address} & \textbf{Network}    \\
    \midrule
    \note{ElectionFactory}           & \note{0x...}     & \note{Polygon Amoy} \\
    \note{Election (implementation)} & \note{0x...}     & \note{Polygon Amoy} \\
    \note{VoterRegistry}             & \note{0x...}     & \note{Polygon Amoy} \\
    \note{VoteForwarder}             & \note{0x...}     & \note{Polygon Amoy} \\
    \bottomrule
  \end{tabular}
\end{table}

\subsection{Divergences from the proposed design}

\stub{List what changed between the Evaluation 1 proposal and the final system,
  with the reason for each. This section is worth writing carefully: a design that
  survived contact with implementation unchanged is usually a design that was not
  tested hard enough.}

% ===========================================================================
\section{Smart Contract Design and Implementation}
% ===========================================================================

\subsection{Contract responsibilities}

\stub{One subsection per contract. For each: its responsibility, its state, its
  entry points, and the access control applied to each entry point.}

\subsection{Election lifecycle}

\stub{Present the lifecycle state machine as implemented, and state which
  transitions are time-gated, which require an explicit call, and which are
  irreversible.}

\subsection{Anonymous eligibility and nullifiers}

\stub{Explain the implemented scheme end to end: how an identity commitment is
  formed, how the off-chain tree is maintained, what the proof attests to, and how
  the nullifier prevents a second ballot without revealing the voter.}

\subsection{Gas-efficiency measures}

\stub{Describe each deliberate gas decision (the off-chain membership tree,
  minimal proxy clones for election deployment, off-chain metadata storage) and
  point forward to the measurement that confirms or refutes each one.}

\subsection{Selected implementation}

\stub{Include the ballot-acceptance function as a listing. Choose one or two
  functions that carry the design's central ideas rather than pasting whole files.}

\begin{lstlisting}[language=Solidity,caption={Ballot acceptance.},label={lst:castvote}]
// Replace with the implemented function.
function castVote(/* ... */) external {
    // verify the membership proof
    // reject a previously seen nullifier
    // record the nullifier and increment the tally
}
\end{lstlisting}

\subsection{Testing and static analysis}

\stub{Report the test suite: unit tests, fuzz campaigns, and what each adversarial
  test asserts. Report the coverage figure and the static-analysis findings by
  severity, with a written disposition for every finding, including the ones
  judged to be false positives, and why.}

% ===========================================================================
\section{Performance Analysis}
% Source material: docs/evaluation-parameters.md, bench/results/
% This is the section the final evaluation weighs most heavily. Every number
% needs a stated environment and a stated sample size.
% ===========================================================================

\subsection{Setup}

\stub{Describe the measurement environment precisely enough to be reproducible:
  hardware, client versions, network, RPC provider, gas price at the time of
  measurement, contract commit hash, sample sizes, and how repetitions were
  aggregated.}

\begin{table}[htbp]
  \centering
  \caption{Measurement environments.}
  \label{tab:bench-setup}
  \small
  \begin{tabular}{@{}L{3.7cm}L{5.3cm}L{6.9cm}@{}}
    \toprule
    \textbf{Environment} & \textbf{Used for} & \textbf{Configuration}                                     \\
    \midrule
    Local EVM            & Gas measurements  & \note{toolchain version, solc version, optimizer settings} \\
    Local node           & Throughput        & \note{host specification}                                  \\
    Public testbed       & Latency, cost     & \note{network, RPC provider, gas price, block time}        \\
    \bottomrule
  \end{tabular}
\end{table}

\stub{State plainly which numbers come from a local environment and which from a
  live network, and why the split was made that way. Do not present a locally
  measured number as though it came from the testbed.}

\subsection{Results by evaluation parameter}

Results are presented against the parameters defined in the Evaluation 1 report.
Targets stated there were hypotheses; where a target was missed, the result is
reported as measured and analysed.

\subsubsection{Cost (A1--A5)}

\stub{Report the gas table and the derived cost per vote. Compare the two
  networks. State whether the minimal-proxy and off-chain-metadata decisions paid
  for themselves, with the measured saving for each.}

\subsubsection{Performance (B1--B5)}

\stub{Report latency distributions rather than averages: median and tail. Give
  proof generation time by device class. Report the throughput curve and say where
  it bends.}

\subsubsection{Scalability (C1--C4)}

This is the central claim of the design.

\begin{figure}[htbp]
  \centering
  \fbox{\parbox[c][5.5cm][c]{0.9\linewidth}{\centering\itshape
      Figure placeholder: gas cost against voter-roll size.\\
      Generate from \texttt{bench/results/} by script, not by hand.}}
  \caption{Gas cost per operation against the size of the electorate.}
  \label{fig:c1}
\end{figure}

\stub{State whether gas is flat in the size of the electorate, and give the
  measured figures at each roll size. If the line is not flat, explain precisely
  where the growth enters and what it would take to remove it. A negative result
  reported clearly is a better outcome than a favourable one that was not measured.}

\subsubsection{Security and correctness (D1--D9)}

\stub{For each property, state how it was demonstrated and what the outcome was.
  Where a property is only partially achieved, say so and give the limit precisely.}

\subsubsection{Usability (E1--E4)}

\stub{Report the user-testing results, including the tasks participants failed
  and what caused each failure. Failures here are the most useful data in the
  section.}

\subsection{Discussion}

\stub{Step back from the individual numbers. Which design decisions were
  vindicated by measurement, which were not, and what would be done differently
  with the benefit of the results? Name the system's principal remaining
  limitation.}

% ===========================================================================
\section{Implementation Plan}
% The template keeps this section in the final report. Use it to record what
% was actually executed against the plan, and what remains.
% ===========================================================================

\stub{Record the plan as executed: what was delivered in each phase, what slipped
  and why, and which contingencies were invoked. Compare against the plan given in
  the Evaluation 1 report rather than presenting a tidied-up version of it.}

\stub{State what remains as future work, with enough specificity to be
  actionable: coercion resistance, tally secrecy while voting is open, and
  production-grade identity verification.}

% ===========================================================================
\section{Conclusion}
% ===========================================================================

\stub{What was built, what the measurements showed, and what the project
  establishes about the viability of the approach. Be measured: state the
  limitations alongside the results.}

% ===========================================================================
\section{References}
% ===========================================================================

\printbibliography[heading=none]

\stub{Cite what was actually read, including anything referenced in the
  Evaluation 1 report that is still relied on here.}

\end{document}
